\documentclass[11pt,a4paper]{article}
\usepackage[T1]{fontenc}
\usepackage[utf8]{inputenc}
\usepackage{amsmath,amssymb,xcolor}
\usepackage[margin=25mm]{geometry}
\usepackage[colorlinks=true,urlcolor=blue]{hyperref}
\definecolor{nexorablue}{HTML}{2459B5}
\setlength{\parindent}{0pt}
\setlength{\parskip}{7pt}
\setlength{\emergencystretch}{2em}
\renewcommand{\familydefault}{\sfdefault}
\begin{document}

{\small\bfseries\color{nexorablue} NEXORA CPGE}\par
{\small MP / NEX-IND-03 / Corrigé}\par
{\LARGE\bfseries Tension d’une bobine en rotation}\par
\vspace{0.25em}\hrule\vspace{0.6em}

\section*{Dérivée du flux}
Le flux lié est $NBS\cos(30t)$. On le dérive et on change le signe. Il vient $e(t)=2{,}16\sin(30t)\,\mathrm V$.

\section*{Valeur efficace}
L’amplitude vaut $2{,}16\,\mathrm V$. La valeur efficace est $2{,}16/\sqrt2\simeq1{,}53\,\mathrm V$. La tension est nulle au moment où le flux est maximal.


\end{document}
